% !TEX root = ../notes.tex

\documentclass[../notes.tex]{subfiles}

\begin{document}

\section{March 16}
Today, we say something about $(\mf g,K)$-modules.

\subsection{\texorpdfstring{$(\mf g,K)$}{(g, K)}-Modules}
The moral of passing to smooth and $K$-finite vectors is that $V^{\mathrm{sm}}_{\mathrm{fin}}$ admits an action by $\mf g$ and an action by $K$.
\begin{lemma} \label{lem:get-gk-mod}
	Fix a continuous representation $V$ of a real Lie group $G$. Then $V^{\mathrm{sm}}_{\mathrm{fin}}$ admits an action by $\mf g$ and an action by $K$.
\end{lemma}
\begin{proof}
	We can see that $K$ sends $K$-finite vectors to $K$-finite vectors, and it also preserves smoothness (because even $G$ sends smooth vectors to smooth vectors). The action by $\mf g$ also turns out to send $K$-finite vectors to $K$-finite vectors: the point is that
	\[kXv=\left(kXk^{-1}\right)kv=(\op{Ad}_kX)(kv).\]
	Thus, $K\cdot\mf g\cdot v=\mf g\cdot K\cdot v$, and $K\cdot v$ is finite-dimensional, so $K\cdot\mf g\cdot v$ is also finite-dimensional.
\end{proof}
\begin{remark}
	On the other hand, $V^{\mathrm{sm}}_{\mathrm{fin}}$ does not, in general, admit an action by $G$ because $G$ does not necessarily send $K$-finite vectors to $K$-finite vectors if $K$ fails to be normal.
\end{remark}
Let's codify the data in the lemma into a definition.
\begin{defihelper}[$(\mf g,K)$-module] \nirindex{Gk module@$(\mf g,K)$-module}
	Fix a real Lie group $G$, let $K\subseteq G$ be the maximal compact subgroup, and let $\mf g$ be the Lie algebra. Then a \textit{$(\mf g,K)$-module} $M$ is a complex vector space which is simultaneously a representation of $\mf g$ and $K$, satisfying the following conditions.
	\begin{listroman}
		\item Locally finite over $K$: for any $m\in M$, the vector space $\op{span}\{km:k\in K\}$ is finite-dimensional.
		\item Differential compatibility: the differential action by $\op{Lie}K$ on $M$ agrees with the action of $\mf g$ on $M$.
		\item Adjoint compatibility: for any $X\in\mf g$ and $k\in K$ and $m\in M$, we have
		\[kXm=(\op{Ad}_kX)kv.\]
	\end{listroman}
	A morphism of $(\mf g,K)$-modules is one preserving the $\mf g$-action and the $K$-action.
\end{defihelper}
\begin{remark}
	On the homework, we will check that (iii) follows from (i) and (ii) if $K$ is connected.
\end{remark}
\begin{example}
	The smooth $K$-finite vectors of a continuous Fr\'echet representation produces a $(\mf g,K)$-module by \Cref{lem:get-gk-mod}, but it is not faithful in general. In particular, it can send non-isomorphic representations of $G$ to isomorphic $(\mf g,K)$-modules. For example, if $G$ acts on a variety $X$, then the representations $L^p(X)$ should give the same $(\mf g,K)$-module for different values of $p$. On the other hand, this functor surjects onto simple modules, and it is an equivalence if one restricts to representations of moderate growth.
\end{example}
\begin{remark}
	The notion of a $(\mf g,K)$-module is rather algebraic: indeed, the $(\mf g,K)$-modules are equivalent to the $(\mf g_\CC,K_\CC)$-modules, and the finite-dimensional continuous representations of $K$ are equivalent to the algebraic representations of $K_\CC$, so everything has been moved to pure algebra. Let's explain this latter equivalence: a choice of finite-dimensional faithful complex representation $V$ of $K$ allows us to define $K_\CC$ as the scheme-theoretic closure of $K\to\op{GL}(V)$. It now turns out that the finite-dimensional representations of $K$ can all be found among the tensor powers of $V$ and $V^*$, so because the representation $K\to\op{GL}(V)$ is tautologically algebraic, algebraicity extends to all representations of $K_\CC$!
\end{remark}
For example, let's classify the $(\mf g,K)$-modules from the group $G=\op{SL}_2(\RR)$. Here, $\mf g=\mf{sl}_2(\RR)$ and $K=\op{SO}_2(\RR)$, which base-changes to $(\mf{sl}_2(\CC),\CC^\times)$.
\begin{example}
	We classify the $(\mf{sl}_2(\CC),\CC^\times)$-modules.
\end{example}
\begin{proof}
	Let $M$ be a $(\mf g_\CC,K_\CC)$-module, and let's start with the $K_\CC$-action. An algebraic representation of $\CC^\times$ amounts to a choice of grading
	\[M=\bigoplus_{n\in\ZZ}M_n,\]
	where $M_n$ is the eigenspace of the $\CC^\times$-action where $\CC^\times$ acts by the $n$th power.
	
	We now turn to the $\mf{sl}_2(\CC)$-action; give $\mf{sl}_2(\CC)$ the standard basis $(e,f,h)$ where $[h,e]=2e$ and $[h,f]=-2f$ and $[e,f]=h$. By changing basis of $\CC^2$, we may assume that $\CC h=\op{Lie}\op{SO}_2(\CC)$. Let's be explicit about this: $\op{SO}(2)$ contains the matrices of the form $\begin{bsmallmatrix}
		\cos\theta & \sin\theta \\ -\sin\theta & \cos\theta
	\end{bsmallmatrix}$, which diagonalizes to $\op{diag}(e^{i\theta},e^{-i\theta})$. We then let $h$ be $i^{-1}\begin{bsmallmatrix}
		& 1 \\ -1
	\end{bsmallmatrix}$ so that $h$ lives in $\mf{sl}_2(\RR)$ (up to the diagonalization!). One also finds that
	\[e=\frac12\begin{bmatrix}
		1 & i \\ i & -1
	\end{bmatrix}\qquad\text{and}\qquad f=\frac12\begin{bmatrix}
		1 & -i \\ -i & -1
	\end{bmatrix}.\]
	Now, by the compatibility of the $\op{Lie}K$-action with the $\mf{sl}_2(\RR)$-action, we see that $h$ must act on $M_n$ by $n$: indeed, $e^{i\theta}$ acts on $M_n$ by $e^{in\theta}$, which differentiates to multiplication by $n$. As for $e$ and $f$, the representation theory of $\mf{sl}_2(\CC)$ shows that the commutator relation $[h,e]=2e$ is equivalent to having $e(M_n)\subseteq M_{n+2}$ for each $n$, and $[h,f]=2f$ means that $f(M_n)\subseteq M_{n-2}$ for each $n$. Lastly, the commutator relation $[e,f]=h$ means that $(ef-fe)|_{M_n}=n\id_{M_n}$.

	One can analyze the simple modules with a little effort, essentially by choosing a basis and then analyzing the possibilities for the $e$-action and $f$-action. Let's just provide the answer.
	\begin{itemize}
		\item Induction: it turns out that most simple $(\mf{sl}_2(\CC),\CC^\times)$-modules thus look like infinite chains of $\CC$s in a single parity, going off infinitely in both directions. These are the ones induced from the subgroup of upper-triangular matrices.
		\item There are the finite-dimensional $(\mf{sl}_2(\CC),\CC^\times)$-modules, which are algebraic and thus all lift to finite-dimensional representations of $\op{SL}_2(\CC)$. The weights look like $\{-m,-m+2,\ldots,m-2,m\}$ for some positive integer $m$.
		\item Discrete series: there are irreducible Verma modules, which start from a given positive integer $m\ge2$ and extend off infinitely in the positive direction; the weights are $m+2\ZZ$. One can also start from a negative integer $m\le-2$ and extend infinitely in the negative direction; the weights are $m-2\ZZ$.
		\item Limits of discrete series: there are irreducible Verma modules starting from $m=1$ and $m=-1$, as in the previous point.
		\qedhere
	\end{itemize}
\end{proof}

\subsection{The Harish-Chandra Isomorphism}
Recall that a representation of a Lie algebra $\mf g$ is equivalent to a representation of the associative algebra $U\mf g$, which is the quotient of the free (non-commutative) algebra $\CC\langle\mf g\rangle$ by the two-sided ideal generated by relations of the form $XY-YX-[X,Y]$.

It is in general difficult to understand non-commutative rings, but sometimes one can try to say something by passing to commutative pieces of it.
\begin{theorem}[Harish-Chandra]
	Fix a finite-dimensional simple complex Lie algebra $\mf g$, and choose a maximal Cartan subalgebra $\mf h\subseteq\mf g$. Then there is an isomorphism
	\[Z(U\mf g)\cong\op{Sym}(\mf h)^W.\]
	Here, the $W$-action on $\mf h$ is $(-\rho)$-shifted.
\end{theorem}
\begin{proof}[Sketch]
	Note that $U\mf g$ admits an adjoint action by the complex algebraic group $G$ corresponding to $\mf g$, defined by extending the adjoint action of $G$ on $\mf g$. Thus, $Z(U\mf g)$ consists exactly of $(U\mf g)^G$.
	
	Now, if we want to produce a $G$-invariant element of $U\mf g$, we could start by noting that $U\mf g$ is a semisimple representation of $G$ (indeed, it is a filtered colimit of finite-dimensional representations), so the surjection $\CC\langle\mf g\rangle\to U\mf g$ descends to a surjection $\CC\langle\mf g\rangle^G\to(U\mf g)^G$. Again, it may be hard to find invariants in $\CC\langle\mf g\rangle$, but we can instead find invariants in the subspace $\op{Sym}\mf g\subseteq\CC\langle\mf g\rangle$. (Note that $\op{Sym}\mf g$ is not a subalgebra!)

	While we're here, let's explicate this construction to find a useful element of $Z(U\mf g)$.
	\begin{definition}[quadratic Casimir]
		Fix a simple Lie algebra $\mf g$. Recall that $\mf g$ admits an invariant non-degenerate bilinear Killing form $K\colon\mf g\otimes\mf g\to\CC$, so there is a canonical identification $\mf g\cong\mf g^*$. We then define the \textit{quadratic Casimir element} $C\in Z(U\mf g)$ to be the image of the identity along the composite
		\[\op{End}(\mf g)=\mf g\otimes\mf g^*\onto\op{Sym}^2\mf g\to U\mf g.\]
	\end{definition}
	\begin{remark}
		Note that $C$ lands in $Z(U\mf g)$ because this composite is $G$-invariant, and $\op{End}(\mf g)^G$ is just $\CC\id_{\mf g}$ because $\mf g$ is simple!
	\end{remark}
	% For example, $\op{Sym}^2\mf g$ contains the $G$-invariant ``quadratic Casimir'' element: r Thus, there is an isomorphism $\mf g\otimes\mf g\cong\mf g\otimes\mf g^*=\op{End}_\CC(\mf g)$ of $G$-representations. The target has a one-dimensional space of $G$-invariants because $\mf g$ is simple, spanned by the identity, which then pulls back to the Casimir element in $\mf g\otimes\mf g$.
	\begin{example}
		For $\mf g=\mf{sl}_2(\CC)$ with standard basis $\{e,f,h\}$, the Killing form is given by $(X,Y)\mapsto\tr(XY)$. One can use this to calculate that the quadratic Casimir element is given by
		\[ef+fe+\frac12h^2.\]
	\end{example}
	Anyway, it turns out that $\op{Sym}(\mf g)^G$ is isomorphic to $Z(U\mf g)$: it is a theorem that $\op{gr}U\mf g\cong\op{Sym}\mf g$, and in particular, one has that our map $\op{Sym}\mf g\to U\mf g$ becomes an isomorphism (of $G$-representations) on the associated graded. It follows that we have an isomorphism of vector spaces $\op{Sym}\mf g\to U\mf g$, and it gives an isomorphism on $G$-invariants before passing to the associated graded by some kind of Five lemma argument.

	We now relate $\op{Sym}(\mf g)^G$ to $\op{Sym}(\mf h)^W$. Well, using the Killing form, we note that there is a natural restriction map
	\[\op{Sym}(\mf g^*)^G\to\op{Sym}(\mf h^*)^W,\]
	which turns out to be an isomorphism by Chevalley's theorem.

	Thus far, we have only provided an isomorphism of vector spaces. It remains to provide an isomorphism of rings. Let $T\subseteq G$ be the maximal torus corresponding to a choice of Cartan $\mf h\subseteq\mf g$. We also choose a polarization so that we have a decomposition $\mf g=\mf n_-\oplus\mf h\oplus\mf n_+$. The moral is that we can start with $Z(U\mf g)=(U\mf g)^G$ and then embed into $(U\mf g)^T$. These invariants come from a grading: indeed, $U\mf g$ can be decomposed as
	\[U\mf g=\bigoplus_{\chi\in X^*(T)}(U\mf g)_\chi,\]
	where $(U\mf g)_\chi$ is the $\chi$-eigenspace. Thus, $(U\mf g)^T$ simply amounts to taking the degree-zero subring. Thus, we may take a quotient by the ideal $I\coloneqq(\mf n_-\cdot U\mf g)_0=(U\mf g\cdot\mf n_+)_0$. (Indeed, one can see that these are the same by pushing all positive or negative root spaces to one side in $U\mf g$.) It turns out that $(U\mf g)_0/I$ is isomorphic to $\op{Sym}(\mf h)^W$ (indeed, we are just killing all monomials which are not purely in $\mf h$), and one checks that this map $Z(U\mf g)\to\op{Sym}(\mf h)^W$ is an isomorphism by comparing it with the isomorphism of vector spaces produced with Chevalley's theorem.
\end{proof}
\begin{example}
	Let's pass the Casimir element through the Harish-Chandra isomorphism. Well, starting with $\frac12h^2+ef+fe$, we first reorder the expression to $\frac12h^2+h+2fe$. Then we take a quotient by $I$, which kills $2fe$, so we get $\frac12h^2+h$ under the Harish-Chandra isomorphism.
\end{example}

\end{document}